\chapter{1993:Kary B. Mullis and Michael Smith}

\label{chap:chemistry1993}

The Nobel Prize in Chemistry for the year 1993 was awarded jointly to Kary B. Mullis and Michael Smith.

\textbf{Motivation:}

for his invention of the polymerase chain reaction (PCR) method (one half, Kary B. Mullis) and for his fundamental contributions to the establishment of oligonucleotide-based, site-directed mutagenesis and its development for protein studies (the other half, Michael Smith)

\section{Kary B. Mullis}

\subsection*{Early Life and Education}

Kary B. Mullis was born in 1944 in Lenoir, North Carolina, USA.

\subsection*{Career and Major Research}

Mullis studied at the Georgia Institute of Technology (B.S., 1966) and earned his Ph.D. in biochemistry from the University of California, Berkeley, in 1973. He joined the Cetus Corporation in 1979.

\subsection*{Nobel Prize-winning Work}

The work for which Kary B. Mullis was awarded the Nobel Prize in Chemistry in 1993 was described by the Nobel Committee as: "for his invention of the polymerase chain reaction (PCR) method".

He invented the polymerase chain reaction, an enzymatic technique that amplifies a specific segment of DNA by making billions of copies automatically.

\subsection*{Significance and Impact}

PCR transformed molecular biology and medicine, enabling DNA cloning, sequencing, diagnosis of diseases, and forensic identification.

\subsection*{Later Career and Honors}

Mullis left Cetus in 1986 and later worked as an independent consultant. He shared the 1993 Nobel Prize in Chemistry, and he died in 2019.

\subsection*{Legacy}

PCR is one of the most widely used techniques in biology.

\subsection*{Selected Publications}

"Specific Enzymatic Amplification of DNA In Vitro: The Polymerase Chain Reaction" (Cold Spring Harbor Symposia on Quantitative Biology, 1986)

\clearpage

\section{Michael Smith}

\subsection*{Early Life and Education}

Michael Smith was born in 1932 in Blackpool, England.

\subsection*{Career and Major Research}

Smith studied at the University of Manchester (B.Sc., 1953; Ph.D., 1956). He emigrated to Canada and joined the University of British Columbia in 1956, where he spent his career.

\subsection*{Nobel Prize-winning Work}

The work for which Michael Smith was awarded the Nobel Prize in Chemistry in 1993 was recognized for his fundamental contributions to the establishment of oligonucleotide-based, site-directed mutagenesis and its development for protein studies.

He developed site-directed mutagenesis, a method to introduce designed mutations at specific positions in a gene using synthetic oligonucleotides, enabling the function of proteins to be studied.

\subsection*{Significance and Impact}

Site-directed mutagenesis became an essential tool in protein engineering and molecular biology.

\subsection*{Later Career and Honors}

Smith remained at the University of British Columbia, where he founded the Biotechnology Laboratory. He shared the 1993 Nobel Prize in Chemistry, and he died in 2000.

\subsection*{Legacy}

The method laid the basis for modern genetic engineering and protein research.

\subsection*{Selected Publications}

"Oligonucleotide-Directed Mutagenesis Using M13-Derived Vectors: An Efficient and General Procedure for the Production of Point Mutations in Any Fragment of DNA" (Nucleic Acids Research, 1982)

\clearpage

\section*{Chapter Summary}

The 1993 Nobel Prize in Chemistry recognized two DNA-based methods that transformed molecular biology. Mullis invented the polymerase chain reaction for amplifying DNA, while Smith developed site-directed mutagenesis for making designed changes in genes.
